% EXTRACTED FOR PAPER TWO (the hardware case study) -- PASS 48.1
% source     : paper/paper.tex
% blob       : 13590dbfe8119022ae95ff7a374910e30146ab8c
% lines      : 902-920 (1-indexed, inclusive)
% section    : Section 5.2
% prose words: 214
%
% The noise-model asymmetry paragraph, including the measured per-qubit readout rates fed back through both routes and the resulting conservative-upper-bound statement. Paper one loses the measured rates and must either cite a rate or drop the quantitative claim.
%
% This is a VERBATIM COPY. The original is unmodified and still frozen.
% Regenerate with: experiments/pass48_extract_paper2.py

The second is the noise model, and it is the stronger assumption. The
channel of Section 4 degrades the collective route only; the single-copy
route is the ideal, noiseless statistical baseline, on the reading that a
joint \(k\)-copy measurement is the harder thing to implement cleanly.
Real shadow readout is not noiseless, and Section 6.3 finds readout to be
the dominant error on the device we tested, so this understates the
single-copy cost, in the direction that costs us, and we have measured
by how much. Simulating readout on both routes at the per-qubit rates
measured in Section 6.3, and correcting the collective parity for readout
exactly through its outcome-dependent per-pair contraction rather than a
blanket factor, moves the crossover to smaller \(n\) in every resolved
cell: by one to two qubits when both routes calibrate readout out, and
further when both leave it uncorrected. The sizes reported here are
therefore a conservative upper bound under this noise model. Under
uncorrected readout the single-copy U-statistic acquires a bias that does
not shrink with the copy budget, unlike its noiseless variance, so both
routes then carry error floors and the comparison changes character. What
the asymmetry does shape is the dependence of \(n^*\) on the noise rate,
which is a property of the collective side by construction.
